\documentclass[11pt,twocolumn]{article}

\usepackage[margin=0.85in]{geometry}
\usepackage[T1]{fontenc}
\usepackage{amsmath,amssymb}
\usepackage{booktabs}
\usepackage{array}
\usepackage{hyperref}
\usepackage{xcolor}
\usepackage{microtype}
\usepackage{authblk}
\usepackage[numbers]{natbib}
\usepackage{enumitem}
\setlist{nosep}

\hypersetup{colorlinks=true,linkcolor=blue!60!black,citecolor=blue!60!black,urlcolor=blue!60!black}

\title{Sex as Phase-Debt Cancellation:\\A Formal Derivation of Recombination Advantage}

\author[1]{Jonathan Washburn}
\affil[1]{Independent Researcher}

\date{March 2026}

\begin{document}
\maketitle

\begin{abstract}
We present a formal model in which sexual reproduction confers a
categorical advantage over asexual reproduction by cancelling accumulated
antisymmetric phase debt. Each chromosome is modeled as an 8-component
phase register, reflecting the 8-tick discrete clock of Recognition
Science. We define the FOLD operator
$\mathrm{FOLD}(v)_k = (v_k + v_{(8-k) \bmod 8})/2$ and prove three
machine-checked theorems in Lean~4: (1)~FOLD is idempotent; (2)~FOLD
annihilates the antisymmetric (imaginary) component of any phase profile;
(3)~asexual reproduction preserves all antisymmetric debt. Together, these
establish that sexual lineages reset their antisymmetric phase debt to
zero each generation, while asexual lineages accumulate debt
monotonically---a formal structural analogue of Muller's ratchet. The
model translates cleanly into the language of standard recombination
theory: the symmetric/antisymmetric decomposition corresponds to
additive/epistatic genetic variance, and FOLD corresponds to the
recombination step that breaks up negative linkage disequilibrium.
\end{abstract}

\section{Introduction}

Why sex exists is one of the oldest open questions in evolutionary
biology. Asexual reproduction is faster, cheaper, and avoids the
``twofold cost of males''~\cite{Maynard1978}. Yet sexual reproduction
is nearly universal among eukaryotes. The classical explanations invoke
Muller's ratchet~\cite{Muller1964} (deleterious mutations accumulate
irreversibly in asexual lineages), the Fisher--Muller
effect~\cite{Fisher1930, Muller1932} (recombination speeds adaptation by
assembling beneficial alleles from different lineages), and
Hill--Robertson interference~\cite{Hill1966} (linked loci interfere with
each other's response to selection).

These explanations are powerful but largely verbal or population-genetic.
They describe \emph{what} sex does (breaks up linkage, restores variation)
but do not identify a single structural operation that captures the
advantage formally. In this paper we provide such an operation.

Recognition Science~\cite{Washburn2025} derives an 8-tick discrete clock
from algebraic first principles. We model each chromosome as an
8-component \textbf{phase register}---a vector
$v \in \mathbb{R}^8$---and define the \textbf{FOLD operator}
\begin{equation}\label{eq:fold}
  \mathrm{FOLD}(v)_k = \frac{v_k + v_{(8-k) \bmod 8}}{2},
  \qquad k \in \{0, \ldots, 7\}.
\end{equation}
We prove three properties of FOLD in Lean~4:
\begin{enumerate}
\item \textbf{Idempotence}: $\mathrm{FOLD}(\mathrm{FOLD}(v)) =
  \mathrm{FOLD}(v)$.
\item \textbf{Antisymmetric annihilation}: the antisymmetric part of
  $\mathrm{FOLD}(v)$ is identically zero.
\item \textbf{Asexual preservation}: asexual reproduction preserves the
  antisymmetric part unchanged.
\end{enumerate}
The biological interpretation is that sexual reproduction physically
instantiates FOLD: meiotic crossover with a complementary partner averages
the two parental phase registers, annihilating accumulated antisymmetric
drift. Asexual lineages skip FOLD, so their antisymmetric component grows
monotonically with each generation of mutation---a phase-profile version
of Muller's ratchet.

\section{Definitions}

\subsection{Phase Profiles}

A \textbf{phase profile} is a function $v : \mathrm{Fin}\,8 \to
\mathbb{R}$. We interpret each component $v_k$ as encoding the
distribution of recognition phases across the chromosome's genomic
content at tick $k$ of the 8-tick cycle.

Every phase profile decomposes uniquely into symmetric and antisymmetric
parts:
\begin{align}
  v^{\mathrm{sym}}_k &= \frac{v_k + v_{(8-k) \bmod 8}}{2},
  \label{eq:sym} \\
  v^{\mathrm{anti}}_k &= \frac{v_k - v_{(8-k) \bmod 8}}{2}.
  \label{eq:anti}
\end{align}
This decomposition is proved to be exact:
$v_k = v^{\mathrm{sym}}_k + v^{\mathrm{anti}}_k$ for all $k$
(\texttt{profile\_decomposition} in
\texttt{ChromosomePhaseProfile.lean}).

The symmetric part represents the ``real'' or additive component of the
phase register. The antisymmetric part represents the ``imaginary'' or
epistatic component---the phase debt that accumulates when the register
drifts away from its conjugate-symmetric ground state.

\subsection{The FOLD Operator}

FOLD is defined by Eq.~\eqref{eq:fold}. Comparing with
Eq.~\eqref{eq:sym}, we see immediately that $\mathrm{FOLD}(v) =
v^{\mathrm{sym}}$: FOLD projects onto the symmetric subspace.

\subsection{Crossover Model}

Meiotic crossover between two parental profiles $p_1, p_2$ is modeled
as pointwise averaging:
\begin{equation}\label{eq:crossover}
  \mathrm{crossover}(p_1, p_2)_k = \frac{(p_1)_k + (p_2)_k}{2}.
\end{equation}
When the second parent has the \textbf{complementary profile}
$p_2 = p_1^{\mathrm{comp}}$, defined by
$(p_1^{\mathrm{comp}})_k = (p_1)_{(8-k) \bmod 8}$, the crossover
reduces to exactly FOLD:
\begin{equation}\label{eq:complement}
  \mathrm{crossover}(p, p^{\mathrm{comp}})_k = \mathrm{FOLD}(p)_k.
\end{equation}
This is proved as \texttt{crossover\_complement\_is\_fold} in Lean.

\section{Results}

\subsection{FOLD Is Idempotent}

\begin{quote}
\textbf{Theorem} (\texttt{fold\_idempotent}). For all
$v : \mathrm{Fin}\,8 \to \mathbb{R}$ and $k \in \mathrm{Fin}\,8$,
$\mathrm{FOLD}(\mathrm{FOLD}(v))_k = \mathrm{FOLD}(v)_k$.
\end{quote}

The proof proceeds by exhaustive case analysis on all 8 values of $k$,
followed by ring normalization. Idempotence means that once a lineage
has applied FOLD (i.e., reproduced sexually once), further applications
do not change the profile. FOLD is a projection.

\subsection{FOLD Annihilates Antisymmetric Phase Debt}

\begin{quote}
\textbf{Theorem} (\texttt{fold\_kills\_antisymmetric}). For all
$p : \mathrm{PhaseProfile}$ and $k \in \mathrm{Fin}\,8$,
$(\mathrm{FOLD}(p))^{\mathrm{anti}}_k = 0$.
\end{quote}

In words: after one application of FOLD, the antisymmetric component of
the resulting profile is identically zero. All accumulated imaginary phase
debt is cancelled. The proof again proceeds by exhaustive
\texttt{fin\_cases} and \texttt{ring}.

\subsection{Asexual Reproduction Preserves All Debt}

\begin{quote}
\textbf{Theorem} (\texttt{asexual\_preserves\_antisymmetric}). For all
$p : \mathrm{PhaseProfile}$ and $k \in \mathrm{Fin}\,8$,
$(\mathrm{asexual}(p))^{\mathrm{anti}}_k = p^{\mathrm{anti}}_k$.
\end{quote}

Asexual reproduction is modeled as the identity on phase profiles: the
offspring inherits the parent's profile unchanged. Since no averaging
occurs, the antisymmetric part is preserved exactly.

\subsection{The Sexual Advantage Theorem}

Combining the three results above:

\begin{quote}
\textbf{Theorem} (\texttt{sexual\_advantage}). The conjunction holds:
\begin{enumerate}
\item For all $p$ and $k$:
  $(\mathrm{FOLD}(p))^{\mathrm{anti}}_k = 0$.
\item For all $p$ and $k$:
  $(\mathrm{asexual}(p))^{\mathrm{anti}}_k = p^{\mathrm{anti}}_k$.
\end{enumerate}
\end{quote}

The interpretation: sexual lineages reset antisymmetric phase debt to zero
each generation. Asexual lineages preserve it. If mutations inject
positive debt $\delta > 0$ per generation, asexual debt grows as
$n\delta$ (proved monotone by \texttt{debt\_monotone}), while sexual
debt is bounded at zero after each mating event. This is a formal
structural analogue of Muller's ratchet.

\section{Translation into Standard Genetics}

The 8-register phase model is abstract, but it maps onto standard
concepts:

\begin{center}
\small
\begin{tabular}{@{}ll@{}}
\toprule
Phase-profile language & Standard genetics \\
\midrule
Phase register $v \in \mathbb{R}^8$ & Genotype \\
Symmetric part $v^{\mathrm{sym}}$ & Additive genetic value \\
Antisymmetric part $v^{\mathrm{anti}}$ & Epistatic / dominance variance \\
FOLD operator & Recombination \\
Complementary profile & Mate with opposing haplotype \\
$\mathrm{FOLD}(v) = v^{\mathrm{sym}}$ & Recombination extracts additive value \\
Asexual $=$ identity & Clonal inheritance \\
Debt accumulation $n\delta$ & Muller's ratchet \\
\bottomrule
\end{tabular}
\end{center}

The key structural insight is that FOLD does not merely ``shuffle''
alleles (as in the Fisher--Muller picture) or ``break up'' linked loci
(as in the Hill--Robertson picture). It performs a specific algebraic
operation---projection onto the symmetric subspace---that has a provably
clean result: the antisymmetric part is exactly zeroed, not merely
reduced.

\subsection{Muller's Ratchet}

In the classical formulation~\cite{Muller1964}, an asexual population
cannot regain a mutation-free genotype once it is lost by drift. In
phase-profile language, the antisymmetric part of the register cannot
decrease without FOLD. Each generation of mutation adds $\delta > 0$ to
the antisymmetric component, and the proof \texttt{debt\_monotone} shows
that accumulated debt is strictly monotone in generation number.

\subsection{Hill--Robertson Interference}

Hill and Robertson~\cite{Hill1966} showed that selection at one locus
interferes with selection at a linked locus. In the phase model, linkage
corresponds to correlation between register components. FOLD decorrelates
the symmetric and antisymmetric parts by projecting onto the symmetric
subspace, eliminating the interference.

\section{Formal Verification}

All results are formalized in Lean~4 (version 4.27.0) with Mathlib. The
two principal modules are:

\begin{itemize}
\item \texttt{TemporalCoordinationLayer.lean}: defines \texttt{foldReal},
  proves \texttt{fold\_idempotent}, and defines the
  \texttt{SexualFold} structure.
\item \texttt{ChromosomePhaseProfile.lean}: defines
  \texttt{PhaseProfile}, \texttt{symmetricPart},
  \texttt{antisymmetricPart}, \texttt{crossoverProfile},
  \texttt{complementaryProfile}, and \texttt{asexualOffspring}. Proves
  \texttt{profile\_decomposition},
  \texttt{crossover\_complement\_is\_fold},
  \texttt{fold\_kills\_antisymmetric},
  \texttt{asexual\_preserves\_antisymmetric},
  \texttt{debt\_monotone}, and the combined
  \texttt{sexual\_advantage}.
\end{itemize}

The formalization contains zero axioms and zero \texttt{sorry} assertions.
Every proof terminates by exhaustive \texttt{fin\_cases} on
$\mathrm{Fin}\,8$ followed by \texttt{ring} normalization, or by direct
computation.

\section{Discussion}

\subsection{What the Model Does and Does Not Show}

The model shows that a single algebraic operation (FOLD on an 8-register)
is sufficient to produce a categorical sexual advantage: bounded vs.\
unbounded antisymmetric debt. This is a theorem about the operator, not
about a particular population model.

The model does \emph{not} show:
\begin{itemize}
\item That chromosomes literally carry 8-component phase registers. This
  is a model-level interpretation.
\item That meiotic crossover literally implements FOLD. The crossover
  model (pointwise averaging with a complementary partner) is a
  simplification; real crossover involves physical strand exchange at
  specific loci.
\item That the twofold cost of males is overcome. The model addresses the
  qualitative advantage of recombination, not the quantitative
  cost--benefit balance.
\end{itemize}

\subsection{Why 8?}

The 8-tick register is not an arbitrary choice. In Recognition Science,
the minimal ledger-compatible walk in three spatial dimensions has period
$2^3 = 8$~\cite{Washburn2025}. The double helix provides $4$ nucleotides
$\times$ $2$ strands $= 8$ recognition slots per position. The FOLD
operator is the natural projection on this 8-periodic structure.

\subsection{Falsifiability}

The model makes two testable predictions:
\begin{enumerate}
\item Asexual lineages should show monotonically increasing epistatic
  variance (phase debt) over evolutionary time, measurable as increasing
  negative linkage disequilibrium.
\item Sexual lineages should show episodic resets of epistatic variance
  correlated with mating events, not gradual decline.
\end{enumerate}

\section{Conclusion}

Sexual reproduction can be understood as the physical instantiation of the
FOLD operator on an 8-component phase register. FOLD projects onto the
symmetric subspace, annihilating accumulated antisymmetric (imaginary)
phase debt. Asexual reproduction preserves all debt. This gives a formal
structural derivation of recombination advantage that is complementary to,
and consistent with, the classical explanations of Muller, Fisher, and
Hill--Robertson. The entire argument is machine-checked in Lean~4.

\section*{Data Availability}

All Lean source files are in the \texttt{reality} repository under
\texttt{IndisputableMonolith/Genetics/}.

\begin{thebibliography}{99}
\bibitem{Fisher1930}
Fisher, R.A. (1930). \textit{The Genetical Theory of Natural Selection}.
Clarendon Press.

\bibitem{Hill1966}
Hill, W.G. and Robertson, A. (1966). The effect of linkage on limits to
artificial selection. \textit{Genet.~Res.} \textbf{8}(3), 269--294.

\bibitem{Maynard1978}
Maynard Smith, J. (1978). \textit{The Evolution of Sex}. Cambridge
University Press.

\bibitem{Muller1932}
Muller, H.J. (1932). Some genetic aspects of sex.
\textit{Am.~Nat.} \textbf{66}(703), 118--138.

\bibitem{Muller1964}
Muller, H.J. (1964). The relation of recombination to mutational advance.
\textit{Mutat.~Res.} \textbf{1}(1), 2--9.

\bibitem{Washburn2025}
Washburn, J. (2025). The Algebra of Reality: A Recognition Science
Derivation of Physical Law. \textit{Axioms} \textbf{15}(2), 90.
\end{thebibliography}

\end{document}
